% ============================================================
% Appendix G: Methodology Detail -- Phase Activities, QA, and Documentation
% ============================================================

\chapter{Methodology: Detailed Phase Activities and Quality Assurance}
\label{app:methodology-detail}

This appendix expands the phase summaries in Chapter~\ref{ch:methodology} with full
activity lists, quality assurance procedures, and documentation and version control practices.

\section{Detailed Phase Activity Lists}

\subsection{Phase 0: Preparation and Infrastructure (July--November 2025)}

\begin{itemize}
\item HomeLab infrastructure initialisation (Proxmox\cite{proxmox}, networking, storage)
\item Development environment setup (VSCode\cite{vscode}, LaTeX\cite{latex_project}, Git)
\item Security tools installation (Parrot OS Security\cite{parrot_os},
      Metasploit\cite{metasploit}, Burp Suite\cite{burp_suite})
\item Technical documentation and baseline establishment
\end{itemize}

\textbf{Deliverables}: Fully operational virtualisation infrastructure; development
environment configuration; baseline documentation of system architecture.

\textbf{Duration}: 105 hours \quad \textbf{Status}: Completed (100\%)

\subsection{Phase 1: Preliminary Project Proposal and Conceptual Design (October 2025--January 2026)}

\begin{enumerate}

\item \textbf{SMART Objectives Definition} (4h) --- Specification of Main Objective (MO):
develop a reusable teaching package. Definition of 4 Specific Objectives (OE1--OE4) with
measurable KPIs. Alignment with institutional GEISI curriculum requirements.

\item \textbf{State of the Art Analysis} (15h) --- Comprehensive review of 7 major
cybersecurity training platforms (HackTheBox, TryHackMe, Cisco Academy, SEED Labs, GOAD,
VulnHub, Hack4u). Identification of pedagogical gaps. Comparative analysis matrix across
multiple dimensions.

\item \textbf{Functional and Non-Functional Requirements} (10h) --- RF1: design and
implement 4 EVE-NG topologies (.unl files). RF2: implement vulnerabilities per
PTES\cite{ptes}, OWASP\cite{owasp_testing_guide}, NIST\cite{nist_framework}, and
CEH\cite{ceh_framework} standards. RF3: create structured assessment materials and rubrics.
NFR1--NFR3 covering performance, framework compliance, and institutional sustainability.

\item \textbf{Pentesting Methodology and Validation Framework} (8h) --- Adaptation of PTES
for educational context. Definition of ethical hacking principles (EC-Council CEH).
Mapping of attack phases to lab objectives. Validation criteria for security effectiveness.

\item \textbf{Design Decisions Justification} (5h) --- Rationale for EVE-NG selection vs
alternative platforms. Technology stack justification (Parrot OS Security, Metasploit, Burp,
OWASP ZAP). Pedagogical approach: hands-on labs with automation.

\item \textbf{Risk Analysis and Mitigation Planning} (15h) --- Identification of 12+ project
risks (technical, pedagogical, institutional). Assessment of impact and probability.
Definition of mitigation strategies.

\item \textbf{Resource Inventory and Allocation} (5h) --- Hardware requirements: dual-socket
server with Proxmox. Software stack: EVE-NG Community Edition, Parrot OS Security images,
security tools. Human resources: student (\textasciitilde{}820h total), faculty supervisor.
Budget considerations and open-source alternatives.

\item \textbf{Gantt Chart and Critical Path Analysis} (20h) --- Development of detailed
project timeline with 37 tasks. Identification of 13 critical path tasks. CSV export and
interactive web-based visualisation. Progress tracking and milestone definition.

\end{enumerate}

\textbf{Deliverables}: Formal project proposal (Preliminary Project Proposal); detailed Gantt chart with
critical paths; risk register and mitigation plan; requirements specification (functional
and non-functional).

\textbf{Duration}: 82 hours \quad \textbf{Status}: Completed \quad
\textbf{Milestone Deadline}: January 16, 2026

\subsection{Phase 2: Implementation and Interim Validation (January--March 2026)}

\begin{enumerate}

\item \textbf{Lab 1: Reconnaissance Development} (14h research + 20h dev) ---
Scenario: intelligence gathering phase using PTES framework.
Tools: Nmap\cite{nmap}, passive reconnaissance techniques.
Topology: multi-tier network with passive monitoring capabilities.
Validation: successful execution of reconnaissance workflow.

\item \textbf{Lab 2: Web Vulnerabilities Development} (14h research + 25h dev) ---
Scenario: OWASP Top 10 vulnerability exploitation.
Tools: Burp Suite, OWASP ZAP, manual testing techniques.
Topology: web application server with vulnerable services.
Validation: exploitation success and reporting accuracy.

\item \textbf{Lab 3: Incident Response Development} (14h research + 25h dev) ---
Scenario: log forensics and incident response (NIST SP 800-61\cite{nist_sp80061} standards).
Tools: Wireshark\cite{wireshark}, tcpdump, network enumeration tools.
Topology: multi-segment network with pre-staged attack artefacts.
Validation: successful timeline reconstruction and flag recovery.

\item \textbf{Pilot Validation with Users} (20h + 6h revision) ---
Informal testing sessions with 4--5 GEISI students conducted as each lab is completed.
Usability observations and timing notes. Post-session survey on satisfaction, difficulty,
and guide clarity. Adjustments applied based on participant feedback.

\item \textbf{Topology Adjustments and Optimisation} (30h) ---
Performance tuning based on pilot feedback.
VM image optimisation and deployment time reduction.
Network configuration refinement for stability.

\item \textbf{Teaching Material Production} (40h) ---
Student lab guide PDF per lab: task description, objectives, hints,
and tool reference. Teacher resolution guide PDF per lab:
one possible walkthrough, flag values, assessment rubric, and common errors.
Technical configuration reference per lab in the repository.

\end{enumerate}

\textbf{Deliverables}: 3 fully functional EVE-NG labs; student and teacher guides;
pilot validation report; technical configuration references.

\textbf{Duration}: 188 hours \quad \textbf{Status}: Completed \quad
\textbf{Milestone Deadline}: March 20, 2026

\subsection{Phase 3: Completion and Final Validation (March--May 2026)}

\begin{enumerate}

\item \textbf{Lab 4: Privilege Escalation Development} (16h research + 30h dev) ---
Scenario: post-exploitation and privilege escalation (CEH framework).
Tools: Metasploit, manual exploitation, system enumeration.
Topology: hardened systems with privilege escalation vectors.
Validation: successful privilege elevation and system compromise.

\item \textbf{Comprehensive Penetration Testing} (30h Round 1 + 25h Round 2) ---
Round 1: sequential testing of Labs 1, 2, and 3.
Round 2: integrated testing across all 4 laboratories.
Documentation of all findings and remediation strategies.

\item \textbf{Automation Scripts Development} (30h) ---
Utility scripts (Bash/Python) for recurring lab management tasks:
bridge configuration, ISO upload, and connectivity check scripts.
Scripts documented and tested in the project EVE-NG environment.

\item \textbf{Final User Validation and Iteration} (18h) ---
Final informal validation sessions with the pilot group (4--5 participants).
Collection of post-session survey data and usability notes.
Implementation of final adjustments based on feedback.

\item \textbf{Bachelor's Thesis Report Writing and Finalisation} (90h) ---
Phase I (15h): Introduction, Objectives, Literature Review.
Phase II (40h): Methodology, Requirements, Implementation Results.
Phase III (35h): Validation, Analysis, Conclusions.
Final revision, bibliography integration, and publication preparation (35h).

\end{enumerate}

\textbf{Deliverables}: 4 fully functional labs; utility automation scripts; final
bachelor's thesis report; pilot validation report; deployment package ready for
institutional use.

\textbf{Duration}: 315 hours \quad \textbf{Status}: In progress \quad
\textbf{Final Milestone Deadline}: May 27, 2026

\section{Quality Assurance Procedures}

\subsection{Technical Validation}

\begin{itemize}
\item Automated testing of lab deployment and reset
\item Manual security testing and penetration verification
\item Performance benchmarking against KPI targets
\item Compatibility testing across EVE-NG versions
\end{itemize}

\subsection{Pedagogical Validation}

\begin{itemize}
\item Alignment with GEISI learning outcomes
\item Assessment of student understanding through rubrics
\item Effectiveness in achieving educational objectives
\item Feedback from faculty and student participants
\end{itemize}

\subsection{Usability Testing}

\begin{itemize}
\item User interface and documentation clarity
\item Accessibility and ease of deployment
\item Support and documentation adequacy
\item User satisfaction scoring ($\geq$~4/5 target)
\end{itemize}

\section{Documentation and Version Control}

\begin{itemize}
\item GitHub repository (\texttt{TheEgea/TFG}) for all code, scripts, and documentation,
      licensed under CC~BY-NC-SA~4.0
\item LaTeX (XeLaTeX + OpenDyslexic) for formal academic documentation (Volume~I and
      Volume~II)
\item MkDocs Material for the operational web documentation site
\item Versioning system for all deliverables with semantic commit messages
\item Change logs and improvement tracking through pull requests
\item Deployment guides and runbooks per lab
\end{itemize}
